\section*{Complete author Section 3.2 main construction}
Now we prove the following Lemma 17 which shows that both $\mathrm{div}$ and $\mathrm{mod}$ have low complexity when the divisor has low-degree and monic (in fact, constant leading-coefficient suffices).
\paragraph{Lemma 17 (Main Lemma).} Let the polynomials $g,h\in\mathbb F[\mathbf x]$ such that $h$ is monic in $x_n$, $L(g)=s_1$, $L(h)=s_2$, and $\deg_{x_n}(h)=d$. Then, both $L(g\mathbin{\mathrm{div}}h),L(g\mathbin{\mathrm{mod}}h)\le O((s_1+s_2)d^2)$.
\begin{proof}Suppose $C$ is a circuit of size $s_2$ which computes $h$ and $C_g$ is the circuit of size $s_1$ which computes $g$. By using Lemma 61, there is a circuit of size $O(s_2d^2)$, which computes $h_0,\ldots,h_{d-1}$. Now, we split every gate $F$ in $C$ into $d+1$ gates $F_0,F_1,\ldots,F_d$. Suppose, the gate $F$ is computing a polynomial $P_F$. Let $P_F\mathbin{\mathrm{mod}}h=\sum_{i<d}p_ix_n^i$. Then we want the property that $F_i$ computes $p_i$ for $i<d$. And if $i=d$ then $F_i$ computes $P_F\mathbin{\mathrm{div}}h$.
Suppose $F$ is a $+$ gate in $C$ with children gates computing the polynomials $a$ and $b$. Again express $a\mathbin{\mathrm{mod}}h=\sum_{i<d}a_ix_n^i$ and $b\mathbin{\mathrm{mod}}h=\sum_{i<d}b_ix_n^i$. It is clear that
\[(a+b)\mathbin{\mathrm{mod}}h=a\mathbin{\mathrm{mod}}h+b\mathbin{\mathrm{mod}}h.\]
Therefore $p_i=a_i+b_i$. It is also clear that $P_F\mathbin{\mathrm{div}}h=a\mathbin{\mathrm{div}}h+b\mathbin{\mathrm{div}}h$.
Suppose $F$ is a $\times$ gate in $C$ with children gates computing the polynomials $a$ and $b$. Again express $a\mathbin{\mathrm{mod}}h=\sum_{i<d}a_ix_n^i$ and $b\mathbin{\mathrm{mod}}h=\sum_{i<d}b_ix_n^i$. It is clear that:
\[(a\cdot b)\mathbin{\mathrm{mod}}h=(a\mathbin{\mathrm{mod}}h\cdot b\mathbin{\mathrm{mod}}h)\mathbin{\mathrm{mod}}h.\]
For $\mathrm{div}$, we have that:
\begin{align*}P_F\mathbin{\mathrm{div}}h&=a\mathbin{\mathrm{div}}h\cdot b\mathbin{\mathrm{div}}h\cdot h+b\mathbin{\mathrm{div}}h\cdot a\mathbin{\mathrm{mod}}h\\&\quad+a\mathbin{\mathrm{div}}h\cdot b\mathbin{\mathrm{mod}}h+(a\mathbin{\mathrm{mod}}h\cdot b\mathbin{\mathrm{mod}}h)\mathbin{\mathrm{div}}h.\end{align*}
We have already computed $a\mathbin{\mathrm{mod}}h,b\mathbin{\mathrm{mod}}h,a\mathbin{\mathrm{div}}h,b\mathbin{\mathrm{div}}h$. So, we only need to compute $(a\mathbin{\mathrm{mod}}h\cdot b\mathbin{\mathrm{mod}}h)\mathbin{\mathrm{mod}}h$ and $(a\mathbin{\mathrm{mod}}h\cdot b\mathbin{\mathrm{mod}}h)\mathbin{\mathrm{div}}h$. Since we have already computed $a_i,b_i$ for all $i<d$, by using Lemma 16, we can compute all the $p_i$ and $(a\mathbin{\mathrm{mod}}h\cdot b\mathbin{\mathrm{mod}}h)\mathbin{\mathrm{div}}h$ in $O((2d-2)d)=O(d^2)$ many gates. Therefore the new circuit has $O(s_1d^2)$ has many gates. Also we used $O(s_2d^2)$ gates to computes $h_0,\ldots,h_{d-1}$. Hence,
\[L(g\mathbin{\mathrm{div}}h)=O((s_1+s_2)d^2),\qquad L(g\mathbin{\mathrm{mod}}h)=O((s_1+s_2)d^2).\]\end{proof}
The following theorem settles Problem 1, when the divisor has small degree (proving Theorem 4).
\paragraph{Proof of Theorem 18.} The above Lemma 17 shows that when $h$ is monic in $x_n$, the upper bound holds. Let $\tau:\mathbb F[\mathbf x]\longrightarrow\mathbb F[\mathbf x]$, be an invertible monic transformation (sends $x_i\mapsto\alpha_i\cdot x_n+x_i$, where $\alpha_i\in\mathbb F$) s.t. $\tau(h)$ is monic wrt $x_n$, such transformation exists (Lemma 68). Note that, $L(\tau(g))\le s+n=O(s_1)$, and $L(\tau(h))\le s_2+n=O(s_2)$. Moreover, as $\tau$ is degree-preserving, $\deg_{x_n}(\tau(h))=d$.
So, apply Lemma 17 to conclude that $\tau(f)=\tau(g)\mathbin{\mathrm{div}}\tau(h)$, has a circuit of size $O((s_1+s_2)d^2)$. We apply $\tau^{-1}$ again (which is just a additive $n$-blowup) to finally deduce that
\[L(f)\le O((s_1+s_2)d^2).\]\qed
\paragraph{Remark 19.} This proof holds when one replaces $L$ by $\tau$, i.e. the constant-free circuit complexity (for definition, see Section 5.2). Note that, neither $\mathrm{div}$ nor $\mathrm{mod}$ introduce any new constant in the process. Moreover, one can choose the $\alpha_i$ to be explicit and $\mathrm{poly}(\log d)$-computable so that $\tau$ is a monic invertible map. This establishes the claim.
\section*{Complete author Section 3.3 border proof}
\paragraph{Proof of Theorem 22.} By definition, there exists $G,H\in\mathbb F(\epsilon)[\mathbf x]$, of size at most $s_1$ and $s_2$, respectively, such that $G:=g+\epsilon\cdot\widetilde g(\mathbf x,\epsilon)$, and $H:=h+\epsilon\cdot\overline h(\mathbf x,\epsilon)$, where $\widetilde g,\overline h\in\mathbb F[\epsilon,\mathbf x]$. We note that, $\deg_{\mathbf x}(H)$ can be larger than $d$. However, using Lemma 62, we know that $L_{\mathbb F(\epsilon)}(\mathrm{Hom}_{\le d}H)\le O(s_2d^2)=:s'_2$.
We denote $\widetilde H:=\mathrm{Hom}_{\le d}H$. It is important to observe that $\widetilde H|_{\epsilon=0}=h$. By definition, there exists $m$ (could be $\exp(s'_2)$) such that
\[\widetilde H:=h+\epsilon\cdot\widetilde h(\mathbf x,\epsilon)=h+\sum_{j\in[m]}\epsilon^jh_j(\mathbf x),\qquad h_j\in\mathbb F[\mathbf x].\]
Let $\tau:\mathbb F[\mathbf x]\longrightarrow\mathbb F[\mathbf x]$, be an invertible monic transformation (sends $x_i\mapsto\alpha_i\cdot x_n+x_i$, where $\alpha_i\in\mathbb F$) s.t. $\tau(h)$ and each $\tau(h_j)$, for $j\in[m]$ is monic wrt $x_n$; such transformation exists (Lemma 68). Note that, $L_{\mathbb F(\epsilon)}(\tau(G))\le O(s_1)$ and $L_{\mathbb F(\epsilon)}(\tau(\widetilde H))\le O(s'_2)$. Further, $\deg_{\mathbf x}(\tau(\widetilde H))=d$, as $\tau$ is a degree-preserving map. We also have the following identities:
\[\tau(\widetilde H)=\tau(h)+\epsilon\cdot\tau(\widetilde h),\qquad\tau(G)=\tau(f)\cdot\tau(h)+\epsilon\cdot\tau(\widetilde g).\]
By assumption, the leading coefficient of $x_n$ in $\tau(\widetilde H)$ (call it $\alpha$) is in $\mathbb F[\epsilon]$ (in fact, $\alpha\not\equiv0\bmod\epsilon$). This basically makes $\tau(\widetilde H)$ a monic polynomial over $\mathbb F(\epsilon)[\mathbf x]$. Therefore, $\mathrm{div}\,\tau(\widetilde H)$ and $\mathrm{mod}\,\tau(\widetilde H)$ now make sense over $\mathbb F(\epsilon)[\mathbf x]$. By simple division, we have
\begin{multline*}\tau(G)\mathbin{\mathrm{div}}\tau(\widetilde H)=\tau(f)\\+\epsilon\cdot(\tau(\widetilde g)-\tau(f)\cdot\tau(\widetilde h))\mathbin{\mathrm{div}}\tau(\widetilde H).\tag{1}\end{multline*}
Note that, Lemma 17 implies $L_{\mathbb F(\epsilon)}(\tau(G)\mathbin{\mathrm{div}}\tau(\widetilde H))=O((s_1+s'_2)d^2)$. By definition of $\underline L$ and Equation (1), it is trivial to conclude that $\underline L(\tau(f))\le O((s_1+s'_2)d^2)=O(s_1d^2+s_2d^4)$. As $\tau$ is invertible, we can get back $f$ by applying $\tau^{-1}$ (incurring $n$-additive blowup). This finally shows
\[\underline L(f)\le O(s_1d^2+s_2d^4).\]\qed
